\subsection{Time Warp: Virtual Time and Recovering from Late Events}
\label{sec:timewarp}

Pcontrol changed the order in which pending computations run. Here we
also give application events simulated times and recover when their
arrival order disagrees with that history.
The \filepath{examples/time-warp/} directory implements a reflective
simulation experiment. A \emph{discrete event simulation} advances by
processing events such as a car arriving or a washer becoming available.
Each event has a simulated time. Processing an event at time 10 does
not require waiting ten real seconds.

Time Warp is an \emph{optimistic} approach: an object processes available
events without first proving that no earlier event will arrive later.
If one does arrive, the object must restore an earlier state and redo
work. This is \emph{rollback}, not an ordinary error retry: the computation
may have been correct for the events known then, but its history was
incomplete. Reflection lets this example put time and recovery policy
in the execution machinery rather than in every application handler.

\subsubsection{Local Time, Receive Time, and Three Queues}

\code{LVT}, the local virtual time, belongs to a simulation object's
metaobject. The custom evaluator makes that value available as the
pseudo-variable \code{LVT} in its scripts. A timed send has the form:

\begin{lstlisting}
[target <= message @ reply :vrt receive-time]
\end{lstlisting}

The evaluator sends the target metaobject a six-field envelope:

\begin{lstlisting}
[:message Message Reply Sender VST VRT]
\end{lstlisting}

\code{VST} is the sender's virtual send time; \code{VRT} is the
intended virtual receive time. Neither is a wall-clock timestamp.
The application still matches \code{Message}, not this entire envelope.
The explicit reply destination remains a separate field.

Distinguish three queues. A simulation metaobject has an ordinary
message queue and an undoable, time-ordered event queue.
Separately, the evaluator's own metaobject has a queue of evaluation
requests. Lowest-LVT-first versus FIFO changes this third queue,
not the timestamp semantics of the event queue.
An ordinary send without \code{:vrt} goes through the ordinary
message path; the source does not automatically timestamp every send.
For example, the washer's self-messages advance its work loop without
advancing virtual time.

\subsubsection{Run the Original Car-Wash Experiment}
\label{sec:tw-loading}

As with pcontrol, start in a scratch copy of the example directory:
\filepath{tw-load.abclr2} uses relative file names and creates
translations and compiled files. The automated check prepares the copy:

\begin{lstlisting}
make -C tutorial timewarp-check
\end{lstlisting}

In a fresh interpreter whose working directory is that copy:

\begin{codeblock}
\begin{lstlisting}
(inc "tw-load" :use-source t)
[starter <== [:start [3 4 6] 3]]
; Runs Lowest LVT first, then FIFO.
; Prints each run's RESULT and counters.
; Replies (:PRIORI (:ROLLBACK ...) (:EXPS ...)
;          :FIFO   (:ROLLBACK ...) (:EXPS ...)).
\end{lstlisting}
\end{codeblock}

\filepath{tw-load.abclr2} first compiles ten files, then incorporates
those same ten files in the same order. Compilation alone does not run
their object definitions (Section~\ref{sec:loading}). Read the loading
sequence as four stages, rather than opening an arbitrary helper first:

\begin{longtable}{L{0.44\textwidth}L{0.46\textwidth}}
\textbf{Files, in loader order} & \textbf{What the stage supplies}\\\hline
\endhead
\filepath{tw-undo-env}, \filepath{tw-undoable-pqueue},
\filepath{tw-output-history} & State frames, replayable input, and
speculative-output history for recovery.\\
\filepath{tw-pqueue}, \filepath{tw-monitor} & Evaluator work queues and
the rollback counter, distinct from the replayable input queue.\\
\filepath{tw-meta}, \filepath{tw-meta-eval}, \filepath{tw-eval} &
Application dispatch/recovery, scheduling of evaluator requests, and
expression evaluation with time/output context.\\
\filepath{tw-group}, \filepath{tw-wash1} & Connect the machinery, then
define the timed car-wash application and its controllers.\\
\end{longtable}

The names in this table omit \code{.abclr2}, just as the loader does.
It does not load the earlier \code{tw-wash0}, the independent
\code{tw-test}, or the unfinished \code{tw-scheduler} fragment.

The service times are 3, 4, and 6 simulated units; the last argument
requests three cars. \code{starter} compares scheduling policies,
whereas \code{starter-1} performs one simulation. Its factories
create a new attendant and new washers for each run.
The returned report contains counters, not the car assignment list.
That list is printed as \code{RESULT=(:FIN (...))}.
Do not assert particular expression counts as a language rule.

\filepath{tw-wash0.abclr2} is an earlier, ordinary-message version:
two queues pair waiting cars with waiting washers, and a two-signal
join completes the starter's request. Its own comment says it has
no rollback or Time Warp; its \code{service-time} argument is not
used to model a timed completion. The loader instead selects
\filepath{tw-wash1.abclr2}, which adds timed events and explicit
simulation completion.

\subsubsection{Follow Three Cars Through the Application}
\label{sec:tw-application}

\code{washers-gen} numbers the washers from zero. It initially
announces washer 0 at virtual time 0, washer 1 at time 1, and washer
2 at time 2. \code{starter-1} injects car arrivals at times 1, 2,
and 3, followed by an arrival sentinel \code{:fin} at time 4.
These initial injections go directly to the attendant's metaobject
with \code{:VRT}; they are not ordinary application messages to
the attendant.

The attendant stores pairs \code{[time car]} and \code{[time washer]}.
When both lists are nonempty, it removes their first entries and
starts washing at the maximum of those two times:

\wholelisting[firstline=79,lastline=87]{examples/time-warp/tw-wash1.abclr2}

This maximum prevents service before either resource is available.
\code{job-record} stores washer, car, and \emph{start} time, not finish
time. Each washer executes a short self-message loop to supply runtime
work, then announces its next availability at
\code{(+ service-time LVT)}:

\wholelisting[firstline=16,lastline=20]{examples/time-warp/tw-wash1.abclr2}

For the checked small run, the assignments are:

\begin{center}
\begin{tabular}{rrrr}
\textbf{Washer} & \textbf{Car} & \textbf{Start} & \textbf{Available again}\\\hline
0 & 1 & 1 & 4\\
1 & 2 & 2 & 6\\
2 & 3 & 3 & 9\\
\end{tabular}
\end{center}

This is a simulated-time trace, not the order in which console lines
must appear. \code{job-record} uses \code{cons}, so the printed
assignment list is in reverse recording order.
After the sentinel, the attendant sets \code{finished} and waits
until all washers are available before sending its result to
\code{starter-1}. The latter joins that result with its own
end-of-injection signal using \code{wait-flag}. The final reply
therefore does not mean merely that all arrival messages were sent.

\subsubsection{What Happens When an Earlier Event Arrives Late?}
\label{sec:tw-recovery}

Open \filepath{tw-meta.abclr2}. The ordinary path resembles the
metaobjects read earlier: enqueue, find a script, evaluate it,
and continue on \code{:end}. The timed path uses \code{:begin-tw}.
For a positive event with \code{VRT >= LVT}, it saves the current
state frame, sets \code{LVT} to \code{VRT}, and evaluates the
selected script with output-history and time added to the context.

Suppose an object has processed time 10 and then receives an event
for time 5. Continuing at 10 would violate the intended event
history. The other branch requests coordinated recovery:

\wholelisting[linerange={96-101,103-105}]{examples/time-warp/tw-meta.abclr2}

\begin{enumerate}
  \item The undoable environment restores a saved frame from before
        the offending receive time and returns its saved time.
  \item The input queue rewinds previously dequeued events at or after
        the cutoff so they can be considered again.
  \item The output history supplies sends made at or after that
        cutoff. \code{anti-messenger} sends matching cancellation
        envelopes to their destinations.
  \item Dispatch resumes, allowing the revised history to be processed.
\end{enumerate}

An \emph{anti-message} is a cancellation record, not an application
request to undo a particular business operation. The queue can remove
a matching positive event that has not yet executed. Cancellation of
already processed work can require the receiving object to roll back,
which can propagate recovery further.
The original queue compares message contents, reply, sender, and times;
it does not introduce a separate globally unique event identifier.
Read these as the experiment's concrete protocols, not a complete
general-purpose distributed simulation framework.

\subsubsection{Read the Three Histories Separately}
\label{sec:tw-histories}

In \filepath{tw-undo-env.abclr2}, \code{:push} saves a frame and
copies its association-list pairs before later binding updates.
\code{backFrame} selects a saved time strictly less than the rollback
cutoff. For a component trace, start with \code{x = 0}, save at 0,
set \code{x = 5}, save at 10, and set \code{x = 8}.
Rollback to 5 restores the frame at 0 and \code{x = 0}.
It does not interpolate a value between snapshots.

This is a shallow copy of binding pairs, not a deep copy of arbitrary
Lisp objects stored in them. Updating a binding through \code{:set}
is different from mutating an array or cons cell reachable through
that binding. Updates delegated to an outer lexical environment also
do not become part of the local frame's saved history.
Do not assume arbitrary host side effects are reversible.

\filepath{tw-undoable-pqueue.abclr2} keeps pending events in
\code{queue} and dequeued events in \code{old}. After processing
events at times 2 and 4, rollback to 3 makes the event at 4 available
again while retaining the earlier history. It also contains
positive/anti-message cancellation and sorting helpers.
Some alternative enqueue code is guarded by \code{\#+COMMENT};
read the active handlers rather than combining all printed alternatives.

\filepath{tw-output-history.abclr2} records the target, message,
reply destination, send time, and receive time. Its cutoff compares the
record's \emph{send time}, the fourth field, not its receive time.
For sends at times 1 and 5, a cutoff of 4 extracts the second send
even if its receive time is 20. Extraction removes those records
from the retained history. Restoring bindings alone would leave
such speculative output effects uncorrected.

\subsubsection{A Complete Two-Object Rollback Cascade}
\label{sec:tw-cascade}

The original three-car run does not force recovery. To see the original
machinery actually restore state, replay inputs, and cancel an output,
use the supplemental application in
\filepath{tutorial/examples/08-timewarp-rollback.abclr2}.
It introduces no new simulation algorithm: the unchanged original
loader, metaobjects, evaluator, environments, and history queues do the
recovery. The new application merely makes their effects easy to inspect.

The sink stores the latest received value and a reverse-order trace:

\wholelisting[firstline=6,lastline=12]{tutorial/examples/08-timewarp-rollback.abclr2}

The source adds each input to its total and sends that total to the sink
one simulated unit later:

\wholelisting[firstline=14,lastline=20]{tutorial/examples/08-timewarp-rollback.abclr2}

Neither application script contains an undo operation. Assignment goes
through the original environment protocol; the timed send is recorded
by the original evaluator. The trace uses \code{cons} to construct a new
list and rebinds \code{trace}, rather than destructively modifying a list
shared with earlier snapshots. That distinction matters for shallow
binding-history restoration.

\paragraph{Arrange a late event deliberately.}
Run \code{make -C tutorial timewarp-check}; the runner copies the
original directory and this supplemental file into
\filepath{build/timewarp-tests/}. For an interactive reproduction,
start a fresh interpreter in that prepared directory and load
\code{(inc "08-timewarp-rollback" :use-source t)}.
The example loads \code{tw-load} and resets the rollback monitor before
creating the two objects. It initially uses Lowest-LVT-first scheduling.

Inject a time-10 event, let the finite computation settle, and only
then inject a time-5 event:

\begin{codeblock}
\begin{lstlisting}
[[meta rollback-source] <= [:message [:add 10] nil nil 0 10]]
; Source total becomes 10 at LVT 10.
; Sink receives (:VALUE 10) at LVT 11.

[[meta rollback-source] <= [:message [:add 1] nil nil 0 5]]
; Source rolls back; its earlier output is cancelled at the sink.
; Final source total = 11, sink value = 11.
\end{lstlisting}
\end{codeblock}

These are deliberate meta-level event injections, not ordinary
application requests. The two \code{nil} fields are the initial reply
destination and sender; the final numbers are \code{VST = 0} and
\code{VRT = 10} or 5. No application reply is requested. The source's
subsequent timed sends name the source object as sender automatically.
The current interactive top-level drains runnable work after each
entered form. The checks exploit that boundary to ensure the first
output has \emph{already executed}, rather than merely waiting in the
sink's queue. This staged arrival order is a test arrangement, not a
promise that arbitrary concurrent clients see the same schedule.

\paragraph{Follow the histories, not just the final sum.}
After the first injection, the source has saved its initial frame at
time 0, set \code{LVT = 10}, and changed \code{total} from 0 to 10.
Its output history contains
\code{[rollback-sink [:value 10] nil 10 11]}.
The sink has likewise saved its initial frame at 0 before processing
time 11. Its value is 10 and its trace is \code{((11 10))}.

Now the late positive event has \code{VRT = 5 < LVT = 10}.
The source's \code{:begin-tw} handler takes the recovery branch.
It restores the saved frame at 0, so both its local time and total
return to zero. Its input queue has already dequeued the offending
time-5 event and the previous time-10 event. Rewinding to 5 makes both
available for replay, sorted by receive time. It must not discard the
late event after using it to trigger recovery.

The source extracts outputs whose send time is at least 5. The old
\code{[:value 10]} send qualifies because its send time was 10;
its receive time 11 is preserved in the cancellation envelope.
Here is the original anti-messenger's active handler:

\wholelisting[firstline=135,lastline=142]{examples/time-warp/tw-meta.abclr2}

The first pattern ends recursion on an empty history with \code{:fin}.
The second uses \code{bracket*} to separate one output record from the
remaining history. It sends directly to the target's metaobject,
changing the envelope's tag to \code{:anti-message}, but preserving
message, reply, sender identity, and both timestamps. Then it processes
the remaining records with the same \code{R}.

The source waits for the anti-messenger's completion before sending
itself \code{:end}. That completion means the cancellation requests
were issued, not that every receiving simulation object has completed
its own rollback. The target sends are asynchronous.
In this checked staged run, the sink receives:

\begin{codeblock}
\begin{lstlisting}
[:anti-message [:value 10] nil rollback-source 10 11]
\end{lstlisting}
\end{codeblock}

Because the positive event has already executed, it is in the sink's
\code{old} input history, not its pending queue. The anti-message
therefore cannot simply remove a pending positive event. The sink
dequeues the anti-message and rolls back to a frame strictly before
11: its initial frame at 0. This restores \code{value = 0} and
\code{trace = nil}. Rewinding the input history brings the old positive
and its anti-message together; \code{remove-anti} removes that pair,
so the invalid \code{[:value 10]} event is not replayed.
The sink has no timed outputs of its own to cancel in this application.

The source then processes its revised inputs at 5 and 10. At 5 it
adds 1 to zero and sends value 1 for receive time 6. At 10 it adds
10 to that corrected state and sends value 11 for receive time 11.
These are new outputs of replayed work, not the old output being
retained under another name. The sink's revised history is therefore:

\begin{center}
\begin{tabular}{rll}
\textbf{Receive time} & \textbf{Application event} & \textbf{Resulting sink value}\\\hline
6 & \code{[:value 1]} & 1\\
11 & \code{[:value 11]} & 11\\
\end{tabular}
\end{center}

The final reverse trace is \code{((11 11) (6 1))}; the speculative
\code{(11 10)} entry has disappeared. The source's final local time is
10, the sink's is 11, and the shared monitor records two rollbacks:
one caused by a late positive input, one caused by cancellation of
already processed output. Console cancellation and rollback lines
remain visible: restoring simulation bindings does not erase printing.

\paragraph{Inspect without injecting another application event.}
The original \code{TW-meta} does not implement an application-state
query. Its own ordinary metaobject can expose the simulation machinery's
bindings. For example:

\begin{codeblock}
\begin{lstlisting}
[source-meta-state := [[meta [meta rollback-source]] <== :state]]
[source-state := [source-meta-state <== [:value-of 'state]]]
[source-state <== [:value-of 'total]] ; => 11
[source-meta-state <== [:value-of 'LVT]] ; => 10
\end{lstlisting}
\end{codeblock}

The inner \code{meta} selects the simulation object's interpreter;
the outer selects that interpreter's ordinary metaobject.
Its \code{:state} reply is an environment whose binding named
\code{state} points to the undoable \emph{application} environment.
\code{LVT}, \code{pqueue}, and \code{output-history} instead belong
to the simulation metaobject's own frame. Do not confuse these two
state environments or treat the returned objects as copied snapshots.

\paragraph{Evidence and limits.}
Twelve assertions run for this cascade with each evaluator scheduler,
in separate fresh processes: Lowest-LVT-first and FIFO. They check
the speculative values and output record, corrected values and trace,
both final local times, the replacement output history, the replayed
source input history, the sink's retained corrected inputs, empty
pending event queues, and the two-rollback counter. Transcript
assertions require the cancellation and both restoration reports.
The finite result is the same under these two checked arrangements;
the test is not a performance comparison.

This closes the gap between isolated history tests and a concrete
metaobject/evaluator recovery cascade. It does not certify every
cancellation arrival order, simultaneous or duplicate event,
zero-time rollback boundary, or timestamped now-type send. In
particular, output-history matching is based on the original event
fields rather than unique event identifiers. The same original
\code{remove-anti} helper compares envelope tails, not the tags;
do not infer a general duplicate-delivery policy from this one pair.
The experimental NR/NS evaluator variants remain separate cases.

\paragraph{Check your understanding.}
Why is restoring only the source's total insufficient? Its speculative
output has already changed the sink. Why is the final value 11 rather
than 1? The time-10 input is replayed after the new time-5 input.
Why does the sink not end with both time-11 values in its trace?
The old input is cancelled during rewind and its earlier state effect
is rolled back. Does waiting for the anti-messenger join the entire
cascade? No; it joins cancellation issuance, not all recipient work.


\subsubsection{The Evaluator and a Source Reading Route}
\label{sec:tw-evaluator}

\filepath{tw-group.abclr2} connects \code{TW-meta-gen} and
\code{TW-eval-gen}. Scheduler replacement transfers the old
evaluator queue's entries to the new queue before installing it.
The queues in \filepath{tw-pqueue.abclr2} support \code{:enq},
\code{:deq}, \code{:listify}, and \code{[:set ...]} for that purpose.
The priority variant sorts by an entry's first field; the normal
variant appends entries for FIFO order.

\filepath{tw-meta-eval.abclr2} wraps evaluation requests as
\code{[LVT Message Reply Sender]} when the message has the
extended eight-field context; other requests receive key zero.
Its counter counts incoming evaluation requests, not cars,
virtual time units, or elapsed execution time.
\filepath{tw-monitor.abclr2} separately counts rollback announcements.

Read the first evaluator definition in \filepath{tw-eval.abclr2}
before its variants. It recognizes compiled expressions, variable
lookup, sends, object/group creation, sequential execution, and
argument-list evaluation. The added \code{Outputs} and \code{LVT}
fields carry the simulation context. A cache indexed by
\code{Id} supplies those fields to subsequent requests lacking
them. Timed past-type sends both construct a timestamped envelope
and record the output; ordinary sends use their separate path.

The same file also defines non-reifying and no-scheduling variants,
selected by the experimental \code{NR} and \code{NS} controllers
in \filepath{tw-wash1.abclr2}. They are not aliases for the two
policies compared by \code{starter}; Section~\ref{sec:tw-variants}
explains their separate setup and evaluator architecture. Timestamped now-type sends need separate
validation; the car-wash application uses timed past-type sends.

\filepath{tw-test.abclr2} supplies independent countdown objects
\code{foo}, \code{bar}, and \code{baz}. Their timed self-sends
increase virtual receive time by 1, 2, and 3 respectively.
Their \code{start} object's handler dispatches the three requests;
it has no final application reply. Do not call it with \code{<==}
and assume it is a completion join.
\filepath{tw-scheduler.abclr2} is an unfinished object-definition
fragment and is not loaded by \code{tw-load}. The operative scheduler
objects come from \code{tw-pqueue}, not that fragment.

\subsubsection{NR and NS: Replacing the Evaluator, Not Just Its Queue}
\label{sec:tw-variants}

The original \code{NR} and \code{NS} controllers in
\filepath{tw-wash1.abclr2} explore different implementations of the
Time Warp evaluator. They are not additional names for the priority
and FIFO policies already compared by \code{starter}.
The latter two replace the \emph{queue} of one schedulable evaluator;
NR and NS install entirely new evaluator objects.

\begin{longtable}{L{0.15\textwidth}L{0.36\textwidth}L{0.39\textwidth}}
\textbf{Variant} & \textbf{Evaluator's metaobject} & \textbf{Meaning of the change}\\\hline
\endhead
Standard & Created by \code{TW-meta-eval-gen}. & Its evaluation-request
queue can be selected as Lowest-LVT-first or FIFO.\\
NR & Explicit \code{non-reifying-meta}. & Evaluation requests use direct,
non-reifying handlers instead of the custom scheduling metaobject.\\
NS & Default ordinary metaobject. & Removes the custom Time Warp
evaluation-request scheduler, but keeps ordinary object dispatch.\\
\end{longtable}

All three evaluator definitions retain a local virtual-time cache and
the extended \code{Outputs}/\code{LVT} context. All are used with
\code{TW-meta-gen} for simulation objects, so selecting NR does not
remove those objects' rollback histories. ``No scheduling'' in NS
does not mean that the runtime stops scheduling processes or ignores
simulation event times. It refers to the omitted customized evaluator
scheduling layer; ordinary queues and the undoable event queue remain.

The group selectors show the replacement directly:

\wholelisting[firstline=22,lastline=31]{examples/time-warp/tw-group.abclr2}

Each creates a fresh evaluator and assigns the group's \code{evaluator}
binding. Unlike \code{:change-scheduler}, these handlers do not transfer
the old evaluator's pending request queue. Nor do they rewrite the
\code{evaluator} binding held by an existing simulation metaobject.
The checks create an object before switching and another afterwards:
the former retains the old evaluator, the latter captures the new one.
This is why the original NR/NS controller switches before
\code{starter-1} creates its attendant and washers. A group's configuration
is a creation policy, not necessarily a live broadcast to all its members.
Use separate fresh sessions for these experiments; restoring the group's
standard evaluator does not retrofit older simulation objects.

\paragraph{Read the controller's result protocol carefully.}
Here is NR's original controller:

\wholelisting[firstline=127,lastline=142]{examples/time-warp/tw-wash1.abclr2}

It selects the new evaluator, resets the rollback monitor, and waits on
the familiar \code{starter-1} simulation. If that request completes,
it prints the result, reads the rollback counter, and replies
\code{[:fifo [:rollback count]]}. NS follows the same structure but sends
\code{:no-schedule-eval}. Both reuse the historical \code{:fifo} result
tag, even though neither installs \code{normal-queue-gen} here. The tag
therefore cannot identify the selected evaluator architecture.
The expression-counter reset/query and \code{:exps} result fields are
commented out. The custom counter belongs to \code{TW-meta-eval}; an
ordinary or non-reifying evaluator metaobject does not acquire its
\code{:count} handler merely because the group has that query method.

The original entry points, after loading \code{tw-load} in a scratch
directory, are:

\begin{codeblock}
\begin{lstlisting}
; Use a fresh session for each variant.
[NR <== [:start [3 4 6] 3]]
; In another fresh session:
[NS <== [:start [3 4 6] 3]]
\end{lstlisting}
\end{codeblock}

Both complete this small case. Use the returned controller result to
identify completion; printed car assignments alone are not a join.

\paragraph{Follow the evaluator's representations.}
The translator represents a simple literal such as \code{:fin} as
\code{[:self-evaluate function]}, where the zero-argument function
returns the literal. Each evaluator invokes that function with
\code{funcall}; a function-valued payload is returned as data.
The standard evaluator shows the protocol directly:

\wholelisting[firstline=43,lastline=44]{examples/time-warp/tw-eval.abclr2}

The \code{[:script scriptObject]} case delegates to
\code{call-script-body}, which accepts internal scripts, object-level
\code{:exec}, or compiled functions. Trace the representation before
deciding which case handles a request.

\paragraph{What the variant checks establish.}
Run \code{make -C tutorial timewarp-variants-check}. Fifteen probes per
variant check a fresh evaluator identity, old/new objects' captured
evaluators, the literal's returned representation and payload,
a synthetic five-argument compiled body that replies 42, and
\code{LVT} lookup at 7. The compiled-body probe exercises
\code{:compiled}, whereas the literal probe exercises
\code{:self-evaluate}; one working case does not certify the other.
Further checks cover numbers, nil, symbols, lists, function-valued payloads,
single invocation of an initializer, and an application script returning 42.

Two car-wash sessions add four checks each: controller result shape,
historical \code{:fifo} tag, and a nonnegative rollback count. Transcripts
also require the \code{:fin} result and all three expected car assignments.
The 53 checks cover these evaluator and application protocols.
Inputs and transcripts stay in \filepath{build/timewarp-variants/}.
Use small runs to understand the control flow before exploring
additional cancellation orders or larger simulations.

\paragraph{Check your understanding.}
Does NR eliminate application-state rollback? No; simulation objects
still use the Time Warp metaobject generator. Does NS mean no queues?
No; it omits a particular scheduling layer. Does a returned \code{:fifo}
tag prove FIFO evaluator scheduling? No; the old controllers reuse that
tag. Can a new group evaluator fix an already created object's evaluator?
Not through these selectors: existing metaobjects retain their captured evaluator.


\paragraph{Check your understanding.}
If a washer starts at time 2 and has service time 4, does its self-loop
advance its virtual time ten times?
\emph{Answer:} No. The ordinary self-messages do not supply timestamps;
the next timed availability is at 6.
Why restore both input and output histories?
\emph{Answer:} Inputs may need replay after state restoration, while
outputs from discarded work may already have affected other objects.

\paragraph{What the checks establish.}
Forty checks run in isolated scratch sessions: twelve exercise
binding restoration, input replay, pending cancellation, and output
history extraction; four check the car-wash report and nonnegative
rollback counters; twenty-four check the two-object rollback cascade
under both evaluator scheduling policies. The transcript additionally checks the three
assignment records under both scheduling policies.
Files stay in \filepath{build/timewarp-tests/}.
The small car-wash case can complete without any rollback, so its
success does not demonstrate a complete rollback cascade.
Section~\ref{sec:tw-cascade} supplies a dedicated recovery integration
check. Other cancellation arrival orders remain outside these recovery
checks; the separate variant probes in Section~\ref{sec:tw-variants}
check small NR/NS runs, not arbitrary rollback schedules.
